% LSPV-1 REPORT::JOINTENDPOINT1::2026-10-12T16:14:02Z
% !TEX program = xelatex
\documentclass[11pt]{article}
\usepackage[a4paper,margin=0.86in]{geometry}
\usepackage{fontspec}
\usepackage{amsmath,amssymb,amsthm,booktabs,array,enumitem,xurl,hyperref,microtype}
\hypersetup{colorlinks=true,linkcolor=blue,urlcolor=blue,citecolor=blue}
\setlength{\emergencystretch}{3em}
\newtheorem{theorem}{Theorem}
\newtheorem{corollary}{Corollary}
\newcommand{\SAT}{\mathsf{SAT}}
\newcommand{\NP}{\mathsf{NP}}
\newcommand{\Hsigned}{\mathcal H^{\pm}_{2,2}}
\newcommand{\EXACT}{\textbf{PROVED/EXACT}}
\newcommand{\KNOWN}{\textbf{KNOWN THEOREM}}
\newcommand{\COMP}{\textbf{COMPUTATIONAL}}
\newcommand{\OPEN}{\textbf{OPEN}}
\newcommand{\RETIRED}{\textbf{FALSIFIED/RETIRED}}
\title{JOINTENDPOINT1: A Time-Dominating Joint Gauge Is the Separation, Not a Weaker Bridge}
\author{LSPV-1 --- Universal Core Lower Bound}
\date{12 October 2026 (UTC)}

\begin{document}
\maketitle

\begin{center}
\fbox{\parbox{0.94\linewidth}{
\textbf{Executive status.} P versus NP remains \OPEN.  On the fixed
NP-complete signed Monotone $3$SAT-$(2,2)$ clean core, the cofinal statement
that every exact uniform deterministic decider has superpolynomial
$T^a(S+1)^b$, for fixed $a>0$ and $b\ge0$, is exactly equivalent to
$P\ne NP$.  The time factor already carries the whole separation; multiplying
by workspace does not create a weaker intermediate obstruction.  The exact
equivalence is \EXACT.  A 21,950-cell rational-exponent and hostile-wrapper
audit passed with no failures and is only \COMP.  The route
\texttt{SUPERPOLY-JOINT-GAUGE-AS-WEAKER-BRIDGE1} is therefore \RETIRED.  A
same-realization landing inside a published \emph{specific} forbidden
time--space region remains \OPEN.}}
\end{center}

\section{Exact core, machine, and quantified target}

Let $K=\Hsigned$ be the duplicate-free proper sign-monochromatic $3$-CNF
family in which every variable occurs exactly twice positively and twice
negatively.  Every clause contains three distinct variables and $3m=4n$.
D\"ocker proves that SAT restricted to $K$ is NP-complete
\cite{Docker20}.  Let $N=|\operatorname{enc}(F)|$ be the ordinary canonical
indexed encoded length.

Fix a standard uniform deterministic read-only-input sequential machine.  For
one total exact decider $A$ of $K$, let $T_A(F)\ge1$ be actual transition count
and $S_A(F)\ge0$ peak writable workspace.  Input, preprocessing, indexing,
simulation, extraction, reconstruction and the one-bit output are all charged.
For fixed rationals $a>0$ and $b\ge0$, define
\[
 G^{a,b}_A(F)=T_A(F)^a\bigl(S_A(F)+1\bigr)^b.
\]
The candidate universal lower bound is the exact cofinal statement
\begin{equation}
 \mathsf{JG}_{a,b}(K):\quad
 \forall A\;\forall d\in\mathbb N\;\forall N_0\in\mathbb N\;\exists F\in K:
 \quad N\ge N_0\ \land\ G^{a,b}_A(F)>N^d.                 \tag{1}
\end{equation}
There is no common hidden clock: $A$ ranges over programs, and $d$ is
quantified after $A$.

\section{Exact endpoint theorem}

\begin{theorem}[\texttt{JOINT-TIME-DOMINATING-ENDPOINT1}; \EXACT]
For every fixed rational $a>0$ and $b\ge0$,
\[
 \mathsf{JG}_{a,b}(K)quad\Longleftrightarrow\quad K\notin\mathsf P
 \quad\Longleftrightarrow\quad P\ne NP.                 \tag{2}
\]
\end{theorem}

\begin{proof}
Suppose $K\in\mathsf P$.  Choose one ordinary uniform decider $A_0$ with
$T_{A_0}(F)\le cN^k$.  On a standard sequential realization, the number of
work cells visited is at most the number of transitions up to a constant, so
$S_{A_0}(F)+1\le c'(N+T_{A_0}(F))\le c''N^{k'}$ for fixed $k'$.  Hence
\[
 G^{a,b}_{A_0}(F)\le C N^{ak+bk'}.
\]
Choosing an integer $d>ak+bk'$ (and enlarging the finite threshold to absorb
$C$) falsifies (1).

Conversely suppose (1) fails.  Then some total exact $A$, some $d$, and some
$N_0$ satisfy
\[
 T_A(F)^a(S_A(F)+1)^b\le N^d
 \quad\text{for all }F\in K\text{ with }N\ge N_0.       \tag{3}
\]
Since $(S_A(F)+1)^b\ge1$, equation (3) gives
$T_A(F)\le N^{d/a}$.  Totality bounds the finitely many encoded inputs below
$N_0$ by a constant.  Thus $A$ is polynomial-time on all of $K$, so
$K\in\mathsf P$.  Finally, the polynomial many-one completeness reduction to
$K$ gives $K\in\mathsf P\Longleftrightarrow P=NP$.
\end{proof}

\begin{corollary}[Polynomially tame monotone gauges]
The same equivalence holds for any nonnegative gauge $Q_N(T,S)$ that (i)
dominates $c^{-1}T^a$ for some fixed $a,c>0$, and (ii) is polynomially bounded
whenever $N,T,S$ are polynomially bounded.  Therefore sums, maxima, and fixed
products containing a positive power of actual time do not yield a weaker
cofinal superpolynomial target.
\end{corollary}

The $a>0$ premise is essential.  RESOURCEAXIS1 already proved that the
analogous workspace-only target is unconditionally false: exhaustive SAT
search uses $O(N)$ writable space.  The present theorem does not silently
upgrade that false statement.

\section{What the theorem falsifies and what it preserves}

The candidate under attack was that replacing a universal time lower bound by
a universal joint time--space product could create an intermediate theorem
strictly weaker than $P\ne NP$.  Equation (2) falsifies that reading whenever
the gauge dominates a positive time power.  This is distinct from the
published FLMV tradeoff.  FLMV excludes selected \emph{polynomial} pairs
$T=N^c,S=N^s$ in a precise random-access model; it does not assert (1), and it
requires both bounds on the same realized machine \cite{FLMV05}.

Thus the useful program remains conditional and local: assuming $P=NP$ gives
some unknown finite polynomial exponent; an independently proved, fully
charged compiler would have to transform that same decider into one specific
FLMV-forbidden pair.  Merely proving (1) proves the desired separation itself.

\section{Finite reproducible audit}

The dependency-free command
\begin{verbatim}
timeout 30s python3 experiments/review/joint_gauge_endpoint_audit.py \
  --output results/raw/review/2026-10-12_joint_gauge_endpoint.json
\end{verbatim}
completed within one second.  It checked five positive rational time powers,
five workspace powers, polynomial time exponents $1$ through $12$, all
workspace exponents at most the time exponent, and target exponents $1$
through $20$.  It also checked 19,200 hostile cells across eight formula
families, four encodings, six wrapper types, and four padding powers.  Totals:
2,250 polynomial-closure checks, 500 nonpolynomial-witness checks, 19,200
hostile cells, and zero failures.

Script SHA-256:
\nolinkurl{2f7657fffb5b893c3f27fd19152fd49c2ad93224627877d0f9b4d297178636a3}.
JSON SHA-256:
\nolinkurl{16cb8d116de370755e5acd5ea112bf2864cb4a4101ab64d27db47166d23c64d3}.
Canonical payload SHA-256:
\nolinkurl{98e70f4baac3bc85f0cf82fcbe52e23bd1b6b8e5859b4e30e444b7b8fbc93f13}.

This audit measures exponent arithmetic only.  It performs no uncharged SAT
preprocessing, compatibility extraction, or reconstruction and makes no
asymptotic inference.  The asymptotic result is the symbolic proof above.

\section{Hostile controls and representations}

Horn, 2SAT, affine/Tseitin Gaussian elimination, and bounded-treewidth dynamic
programming are polynomial easy islands; independent products can be streamed.
They show again that semantic multiplicity and compatibility structure do not
force large memory.  They do not refute (2), because a counterexample to (1)
must be one polynomial algorithm on the entire NP-hard core.  Padding,
redundant clauses, variable permutations, complementations, auxiliary-variable
encodings, compact circuits and branching programs change the representation
or restricted model.  Once actual encoded length and conversion costs are
charged, a fixed polynomial presentation remains polynomial; none supplies a
superpolynomial lower bound for arbitrary computation.

DEAD, COPY, DELAY, INTERP and RECOMP wrappers may inflate exact exponents, as
JOINTGAUGE1 proves.  They cannot turn an existing polynomial decider into a
superpolynomial one when wrapper powers are fixed.  Conversely, quotienting
programs by extensional behavior does not create an effective canonical
minimum.  The endpoint theorem is presentation-robust precisely because it
asks only whether \emph{some} polynomial decider exists.

A one-bit decision output need not expose a witness, compatibility table,
proof, answer vector, or chosen protocol.  The theorem uses only charged
machine time and space and therefore supplies no solver-to-protocol transfer.
Nonuniform circuits are outside its quantifiers; the corresponding unrestricted
circuit target is $NP\not\subseteq P/poly$, as CIRCEQUIV1 records.

\section{Barrier and cross-domain transfer audit}

The proof of (2) relativizes as a schema for an oracle-relative complete core;
it therefore does not evade Baker--Gill--Solovay \cite{BGS75}.  It uses no
nonalgebrizing step and supplies no large constructive circuit property;
Aaronson--Wigderson and Razborov--Rudich remain diagnostics, not positive
inferences \cite{AW09,RR97}.  The hidden premise audit is exact: assuming a
universal superpolynomial time-dominating gauge is assuming the separation in
equivalent form.

\begin{center}\small
\begin{tabular}{@{}p{0.25\linewidth}p{0.25\linewidth}p{0.42\linewidth}@{}}
\toprule
Source domain & Classification & Exact correspondence or missing theorem\\ \midrule
FLMV SAT time--space & Direct theorem transfer & Only after one ordinary random-access SAT realization meets the stated simultaneous bounds.\\
Borodin--Cook sorting and Thompson VLSI & Useful analogy requiring a new transfer & Their outputs, access models, layout, routing and area measures are not arbitrary one-bit SAT computation.\\
Communication, pebbling, HEC/RCCN & Useful analogy requiring a new transfer & Need a static-SAT compiler preserving cuts, rereads, recomputation and the same realization.\\
Generic ``resource product'' rhetoric & Misleading / failed transfer & A cofinal time-dominating product is exactly the separation; a restricted product theorem cannot be imported by name.\\
\bottomrule
\end{tabular}
\end{center}

\section{Post-quantum technique ledger}

Schur squares, shortening/puncturing, quadratic-relation ranks, Hilbert
functions, syzygies and graded Betti numbers are \textbf{FAILED TRANSFERS} for
(1) without a theorem forcing every decider into that representation; hidden
algebraic structure may instead enable an algorithmic shortcut.  Collision
entropy and $|\mathrm{im},L|\ge |A|^2/\mathrm{Coll}(L)$ are
\textbf{EXACT OBJECT/THEOREM TRANSFERS} only for a declared map $L$; no such
map is forced here.  LWE-style worst-case/average-case extraction and
special-soundness/collision extraction are \textbf{ANALOGY/CANDIDATE} routes
needing an unconditional black-box bridge.  Cryptographic binding, hashing,
and random-oracle assumptions are \textbf{FAILED TRANSFERS} to an
unconditional P-vs-NP proof.

\section{CORPUS REVIEWED}

The current GitHub tree was enumerated at 2,628 paths.  The run read root
\path{main.tex}, \path{AGENTS.md}, \path{CODEX.md}, the complete
pnp-research skill and protocol references, actual \path{main4.tex} and
\path{main5.tex}, and every current task specification.  It verified newest
MAIN \path{research/main/2026-10-12_1547_MAIN_SYNTHESIS-0025.tex} and newest
LSPV-1 through LSPV-4 reports DERIVEDLEAK1, PRODSTREAM1, FIELDMASK1, and
MODELBRIDGE-0028.  Material historical reports included RESOURCEAXIS1,
HEVALLB1, KCALIB1, TDLG1, MAJORANTNOGO1, POINTINF1, CIRCEQUIV1, HALTCANON1,
CUTGAP1, JOINTGAUGE1 and COREEQUIV1.  Repository search covered joint gauges,
time--space products, cofinal lower bounds, arbitrary exponents, padding and
resource-axis equivalences.

The underlying HEC, VSC-HEC, G-PPC-HEC, CRS-PPC, CCSPC and RCCN causal
components were inspected, together with governance and the current tasks.
Newest MAIN has the required header and final sentinel, so predecessor recovery
was unnecessary.  Primary literature was checked before the substantive
claim: D\"ocker, FLMV, Baker--Gill--Solovay, Aaronson--Wigderson,
Razborov--Rudich, Borodin--Cook, and Thompson.

\section{DUPLICATION / CONTRADICTION AUDIT}

RESOURCEAXIS1 already proved the time-coordinate biconditional and the false
workspace-only analogue.  JOINTGAUGE1 proved joint padding dependence of
exact exponents.  The new contribution is the quantified closure of the gap
between them: every fixed joint gauge that dominates a positive time power has
the same cofinal endpoint as time alone.  This is a genuine extension, not a
renaming of either prior theorem.

There is no contradiction with FLMV, Williams-type product endpoints, or
restricted sorting/VLSI tradeoffs.  Those results exclude particular
polynomial regions in particular models.  Equation (2) classifies the much
stronger all-polynomial endpoint.  There is no contradiction with the linear
space exhaustive algorithm because $a>0$ is essential.

\section{WHAT CHANGED RELATIVE TO ROOT MAIN}

Root MAIN already records the time-only equivalence and joint wrapper gauge.
It should now state their exact synthesis: multiplying by workspace does not
weaken a cofinal superpolynomial target when the gauge retains a positive time
power.  Therefore the next lower-bound attempt must land in a selected
published forbidden pair, not seek a generic superpolynomial joint gauge.
No main4/main5 construction claim is strengthened; ARC remains an achievable-
region program and WRAPPER remains a partial construction program.

\section{ADOPTED / REJECTED / RETIRED PRIOR COMPONENTS}

\textbf{Adopted:} D\"ocker's hard core, RESOURCEAXIS1, COREEQUIV1,
JOINTGAUGE1's wrapper warning, FLMV only in its matched model,
SAME\_REALIZATION, and full cost charging.

\textbf{Rejected:} semantic-state count as runtime; workspace multiplication
as an automatic weakening; compact circuit/branching-program representations
as lower bounds for arbitrary uniform machines; and external validator or
protocol cost charged to the original decider.

\textbf{FALSIFIED/RETIRED:}
\texttt{SUPERPOLY-JOINT-GAUGE-AS-WEAKER-BRIDGE1}.  Not retired are published
restricted-model tradeoffs, calibrated polynomial-region landings, or an
endpoint-free nonrelativizing law if one is actually proved.

\section{STALE OR OVERSTATED CLAIMS FOUND}

Any claim that a universal superpolynomial $T(S+1)$ or $T^a(S+1)^b$ theorem is
an easier intermediate than $P\ne NP$ is overstated on an NP-complete clean
core.  Any claim that a time--space product theorem transfers merely because
another model also names ``time'' and ``space'' is model-mismatched.  Any
claim suppressing $a>0$, encoded length, finite exceptional inputs, or the
order $\forall A\forall d\forall N_0\exists F$ is under-specified.

\section{Status ledger and handoff}

\begin{center}\small
\begin{tabular}{@{}p{0.50\linewidth}p{0.20\linewidth}p{0.22\linewidth}@{}}
\toprule
Item & Status & Scope\\ \midrule
\texttt{JOINT-TIME-DOMINATING-ENDPOINT1} & \EXACT & Fixed $a>0,b\ge0$ on $K$\\
Signed-core NP-completeness & \KNOWN & D\"ocker family\\
Finite exponent/wrapper audit & \COMP & 21,950 checks\\
\texttt{SUPERPOLY-JOINT-GAUGE-AS-WEAKER-BRIDGE1} & \RETIRED & Equivalent to separation\\
Selected same-realization FLMV landing & \OPEN & Fully charged compiler\\
Nonrelativizing/nonalgebrizing hinge & \OPEN & Arbitrary computation\\
$P\ne NP$ & \OPEN & No proof\\
\bottomrule
\end{tabular}
\end{center}

The highest-value lower-bound target remains: from an arbitrary assumed
polynomial SAT decider with unknown finite exponent, construct on the same
ordinary realization a fully charged solver inside one cited unconditional
forbidden time--space region, surviving Horn, affine/Tseitin, 2SAT,
bounded-treewidth, padding, wrapper, one-bit-output and alternate-
representation controls.  This arrow is \OPEN.

\begin{thebibliography}{9}
\bibitem{Docker20} J. D\"ocker, ``Monotone 3-SAT-(2,2) is NP-complete,''
\emph{Theoretical Computer Science} 812 (2020), 1--9,
\url{https://arxiv.org/abs/1912.08032}.
\bibitem{FLMV05} L. Fortnow, R. Lipton, D. van Melkebeek and A. Viglas,
``Time-Space Lower Bounds for Satisfiability,'' \emph{JACM} 52(6), 2005,
\url{https://pages.cs.wisc.edu/~dieter/Research/sat-ts.html}.
\bibitem{BGS75} T. Baker, J. Gill and R. Solovay, ``Relativizations of the
$P=?NP$ Question,'' \emph{SIAM Journal on Computing} 4(4), 1975,
\url{https://doi.org/10.1137/0204037}.
\bibitem{AW09} S. Aaronson and A. Wigderson, ``Algebrization: A New Barrier
in Complexity Theory,'' \emph{ACM TOCT} 1(1), 2009,
\url{https://eccc.weizmann.ac.il/report/2008/005/}.
\bibitem{RR97} A. Razborov and S. Rudich, ``Natural Proofs,'' \emph{JCSS}
55(1), 1997, \url{https://eccc.weizmann.ac.il/eccc-reports/1994/TR94-010/}.
\bibitem{BC82} A. Borodin and S. Cook, ``A Time-Space Tradeoff for Sorting on
a General Sequential Model of Computation,'' \emph{SIAM J. Comput.} 11(2),
1982, \url{https://doi.org/10.1137/0211022}.
\bibitem{Tho79} C. Thompson, ``Area-Time Complexity for VLSI,'' STOC 1979,
\url{https://doi.org/10.1145/800135.804401}.
\end{thebibliography}

\noindent\textbf{Final status:} P versus NP remains \OPEN.

% END_OF_VERIFIED_REPORT::LSPV-1::JOINTENDPOINT1
\end{document}
